\scanpage{64}
\mspage{31}

conclut que $\operatorname{Gl}_\Lambda(M)$ opère par l'intermédiaire d'une \deleted{opération}

\[
 \begin{aligned}
 &\text{opération de } \inserted{G_M^{\deleted{\unsure{*}}}=}\operatorname{Gl}_\Lambda(M)' \inserted{=\operatorname{Gl}_\Lambda(M)/{\pm1}}\text{ sur } S_M,\\
 &\text{respectant la structure } A_M\subset\unsure{\mathscr P|\mathcal P|\mathfrak P}_2(S_M),\\
 &\text{et l'ens. quotient }\Pi_M\text{ de }A_M,
 \end{aligned}\tag*{(164)}
\]

% slanted
\margin{\uline{NB} l'opération de $G_M$ sur $\Pi_M$ respecte aussi sa structure de $Z'_\Lambda$-torseur, \deleted{\unsure{où}} $Z'_\Lambda=\Lambda^*/{\pm1_\Lambda}$, et $G_M$ opère sur $\Pi_M$ via $\delta_M$.}

pour $i\in\Pi$, $g\in G_M$, $g_{S_M}$ est un isom. de la structure $(S_M,A_{M,i})$ sur la structure $(S_M,A_{M,g_{\Pi_M}i})$. En particulier, si on \deleted{\unsure{pose}} considère le composé

\diagram{\begin{tikzcd}[ampersand replacement=\&]
 \text{(165)} \& G_M \arrow[r,"\text{det}"] \arrow[rr,bend right=15,"\delta_M"'] \& \Lambda^* \arrow[r,"\text{can}"] \& \Lambda^*/{\pm1_\Lambda}\overset{\text{def}}{=}Z'_\Lambda{,}
\end{tikzcd}}

le stabilisateur \inserted{dans $G_M$} de \deleted{tout} point de $\Pi_M$ \deleted{est} le noyau \inserted{$G_M^0={}$} $\operatorname{Ker}\delta_M$ intervenant dans la suite exacte

\diagram{\begin{tikzcd}[ampersand replacement=\&]
 \text{(\unsure{166|165})} \& 1\arrow[r,-] \& G_M^0\arrow[r] \& G_M\arrow[r,"\delta_M"] \& Z'_\Lambda\arrow[r,-] \& 1{,}\\
 \& \& {\unsure{\operatorname{Sl}|\operatorname{Gl}}}_\Lambda^1(M)^{\unsure{\pm1|/{\pm1}}}\arrow[u,equal] \& \operatorname{Gl}_\Lambda^{\unsure{*}}(M)'\arrow[u,equal]
\end{tikzcd}}

et il en résulte que c'est aussi le sous-groupe de $G_M$ qui stabilise la structure de \deleted{graphe} graphe $\Sigma(i)$, pour un $i\in\Pi_M$ quelconque donné.

% archive sheet number: 63
